\section{Spectral sequences: filtered complexes}
\label{section-filtered-complex}

\begin{definition}
\label{definition-filtered-complex}
Let $\mathcal{A}$ be an abelian category.
A {\it filtered complex $K^\bullet$ of $\mathcal{A}$}
is a complex of $\text{Fil}(\mathcal{A})$ (see
Definition \ref{definition-filtered}).
\end{definition}

\noindent
We will denote the filtration on the objects by $F$. Thus
$F^pK^n$ denotes the $p$th step in the filtration of the $n$th term of the
complex. Note that each $F^pK^\bullet$ is a complex of $\mathcal{A}$.
Hence we could also have defined a filtered complex as a filtered object
in the (abelian) category of complexes of $\mathcal{A}$.
In particular $\text{gr} K^\bullet$ is a graded object of the
category of complexes of $\mathcal{A}$.

\medskip\noindent
To describe the spectral sequence associated to such an object
we assume, for the moment, that $\mathcal{A}$ is an abelian category
which has countable direct sums and countable direct sums are exact
(this is not automatic, see Remark \ref{remark-direct-sums-not-exact}).
Let us denote $d$ the differential of $K$. Forgetting the grading
we can think of $\bigoplus K^n$ as a filtered differential object of
$\mathcal{A}$. Hence according to Section \ref{section-filtered-differential}
we obtain a spectral sequence $(E_r, d_r)_{r \geq 0}$. In this section
we work out the terms of this spectral sequence, and we endow the terms
of this spectral sequence with additional structure coming from the
grading of $K$.

\medskip\noindent
First we point out that $E_0^p = \text{gr}^p K^\bullet$ is a
complex and hence is graded. Thus $E_0$ is bigraded in a natural
way. It is customary to use the bigrading
$$
E_0 = \bigoplus\nolimits_{p, q} E_0^{p, q},
\quad
E_0^{p, q} = \text{gr}^p K^{p + q}
$$
The idea is that $p + q$ should be thought of as the {\it total degree} of
the (co)homology classes. Also, $p$ is called the {\it filtration degree},
and $q$ is called the {\it complementary degree}.
The differential $d_0$ is compatible with this
bigrading in the following way
$$
d_0  = \bigoplus d_0^{p, q},
\quad
d_0^{p, q} : E_0^{p, q} \to E_0^{p, q + 1}.
$$
Namely, $d_0^p$ is just the differential on the complex
$\text{gr}^p K^\bullet$ (which occurs as $\text{gr}^pE_0$ just shifted
a bit).

\medskip\noindent
To go further we identify the objects $B_r^p$ and $Z_r^p$ introduced
in Section \ref{section-filtered-differential} as graded objects and
we work out the corresponding decompositions of the differentials.
We do this in a completely straightforward manner, but again we warn
the reader that our notation is not the same as notation found
elsewhere. We define
$$
Z_r^{p, q} =
\frac{F^pK^{p + q} \cap d^{-1}(F^{p + r}K^{p + q + 1}) + F^{p + 1}K^{p + q}}
{F^{p + 1}K^{p + q}}
$$
and
$$
B_r^{p, q} =
\frac{F^pK^{p + q} \cap d(F^{p - r + 1}K^{p + q - 1}) + F^{p + 1}K^{p + q}}
{F^{p + 1}K^{p + q}}
$$
and of course $E_r^{p, q} = Z_r^{p, q}/B_r^{p, q}$.
With these definitions it is completely clear that
$Z_r^p = \bigoplus_q Z_r^{p, q}$,
$B_r^p = \bigoplus_q B_r^{p, q}$, and
$E_r^p = \bigoplus_q E_r^{p, q}$. Moreover, we have
$$
0 \subset \ldots \subset B_r^{p, q} \subset
\ldots
\subset Z_r^{p, q} \subset \ldots \subset E_0^{p, q}
$$
Also, the map $d_r^p$ decomposes as the direct sum of the maps
$$
d_r^{p, q} : E_r^{p, q} \longrightarrow E_r^{p + r, q - r + 1},
\quad
z + F^{p + 1}K^{p + q}
\mapsto
dz + F^{p + r + 1}K^{p + q + 1}
$$
where $z \in F^pK^{p + q} \cap d^{-1}(F^{p + r}K^{p + q + 1})$.

\begin{lemma}
\label{lemma-spectral-sequence-filtered-complex}
Let $\mathcal{A}$ be an abelian category. Let $(K^\bullet, F)$ be a
filtered complex of $\mathcal{A}$. There is a spectral sequence
$(E_r, d_r)_{r \geq 0}$ in the category of bigraded objects of $\mathcal{A}$
associated to $(K^\bullet, F)$ such that $d_r$ has bidegree $(r, - r + 1)$
and such that $E_r$ has bigraded pieces $E_r^{p, q}$ and maps
$d_r^{p, q} : E_r^{p, q} \to E_r^{p + r, q - r + 1}$ as given above.
Furthermore, we have $E_0^{p, q} = \text{gr}^p(K^{p + q})$,
$d_0^{p, q} = \text{gr}^p(d^{p + q})$,
and $E_1^{p, q} = H^{p + q}(\text{gr}^p(K^\bullet))$.
\end{lemma}

\begin{proof}
If $\mathcal{A}$ has countable direct sums and if countable direct
sums are exact, then this follows from the discussion above.
In general, we proceed as follows; we strongly suggest the reader
skip this proof. Consider the bigraded object $A = (F^{p + 1}K^{p + 1 + q})$
of $\mathcal{A}$, i.e., we put $F^{p + 1}K^{p + 1 + q}$ in degree $(p, q)$
(the funny shift in numbering to get numbering correct later on).
We endow it with a differential $d : A \to A[0, 1]$ by using $d$
on each component. Then $(A, d)$ is a differential bigraded object.
Consider the map
$$
\alpha : A \to A[-1, 1]
$$
which is given in degree $(p, q)$ by the inclusion
$F^{p + 1}K^{p + 1 + q} \to F^pK^{p + 1 + q}$.
This is an injective morphism of differential objects
$\alpha : (A, d) \to (A, d)[-1, 1]$. Hence, we can apply
Remark \ref{remark-differential-object-selfmap}
with $S = [0, 1]$ and $T = [1, -1]$.
The corresponding spectral sequence $(E_r, d_r)_{r \geq 0}$
of bigraded objects is the spectral sequence we are looking
for. Let us unwind the definitions a bit. First of all we have
$E_r = (E_r^{p, q})$. Then, since $T^rS = [r, -r + 1]$
we have $d_r : E_r \to E_r[r, -r + 1]$ which means that
$d_r^{p, q} : E_r^{p, q} \to E_r^{p + r, q - r + 1}$.

\medskip\noindent
To see that the description of the graded pieces hold, we argue
as above. Namely, first we have
$$
E_0 = \Coker(\alpha : A \to A[-1, 1])[0, -1] =
\Coker(\alpha[0, -1] : A[0, -1] \to A[-1, 0])
$$
and by our choice of numbering above this gives
$$
E_0^{p, q} = \Coker(F^{p + 1}K^{p + q} \to F^pK^{p + q}) = \text{gr}^pK^{p + q}
$$
The first differential is given by
$d_0^{p, q} = \text{gr}^pd^{p + q} : E_0^{p, q} \to E_0^{p, q + 1}$.
Next, the description of the boundaries $B_r$ and the cocycles $Z_r$
in Remark \ref{remark-differential-object-selfmap}
translates into a straightforward manner into the formulae
for $Z_r^{p, q}$ and $B_r^{p, q}$ given above.
\end{proof}

\begin{lemma}
\label{lemma-spectral-sequence-filtered-complex-d1}
Let $\mathcal{A}$ be an abelian category.
Let $(K^\bullet, F)$ be a filtered complex of $\mathcal{A}$.
Assume $\mathcal{A}$ has countable direct sums.
Let $(E_r, d_r)_{r \geq 0}$ be the spectral sequence
associated to $(K^\bullet, F)$.
\begin{enumerate}
\item The map
$$
d_1^{p, q} :
E_1^{p, q} = H^{p + q}(\text{gr}^p(K^\bullet))
\longrightarrow
E_1^{p + 1, q} = H^{p + q + 1}(\text{gr}^{p + 1}(K^\bullet))
$$
is equal to the boundary map in cohomology associated to the short
exact sequence of complexes
$$
0 \to \text{gr}^{p + 1}(K^\bullet) \to
F^pK^\bullet/F^{p + 2}K^\bullet \to \text{gr}^p(K^\bullet) \to 0.
$$
\item Assume that $d(F^pK) \subset F^{p + 1}K$ for all $p \in \mathbf{Z}$.
Then $d$ induces the zero differential on $\text{gr}^p(K^\bullet)$
and hence
$E_1^{p, q} = \text{gr}^p(K^\bullet)^{p + q}$.
Furthermore, in this case
$$
d_1^{p, q} :
E_1^{p, q} = \text{gr}^p(K^\bullet)^{p + q}
\longrightarrow
E_1^{p + 1, q} = \text{gr}^{p + 1}(K^\bullet)^{p + q + 1}
$$
is the morphism induced by $d$.
\end{enumerate}
\end{lemma}

\begin{proof}
This is clear from the formula given for the differential
$d_1^{p, q}$ just above Lemma \ref{lemma-spectral-sequence-filtered-complex}.
\end{proof}

\begin{lemma}
\label{lemma-spectral-sequence-filtered-complex-functorial}
Let $\mathcal{A}$ be an abelian category.
Let $\alpha : (K^\bullet, F) \to (L^\bullet, F)$ be a morphism of
filtered complexes of $\mathcal{A}$. Let $(E_r(K), d_r)_{r \geq 0}$,
resp.\ $(E_r(L), d_r)_{r \geq 0}$ be the spectral sequence associated
to $(K^\bullet, F)$, resp.\ $(L^\bullet, F)$.
The morphism $\alpha$ induces a canonical morphism of spectral
sequences $\{\alpha_r : E_r(K) \to E_r(L)\}_{r \geq 0}$ compatible
with the bigradings.
\end{lemma}

\begin{proof}
Obvious from the explicit representation of the terms of the
spectral sequences.
\end{proof}

\begin{definition}
\label{definition-filtration-cohomology-filtered-complex}
Let $\mathcal{A}$ be an abelian category.
Let $(K^\bullet, F)$ be a filtered complex of $\mathcal{A}$.
The {\it induced filtration} on $H^n(K^\bullet)$ is the filtration defined
by $F^pH^n(K^\bullet) = \Im(H^n(F^pK^\bullet) \to H^n(K^\bullet))$.
\end{definition}

\noindent
Writing out what this means we see that
\begin{equation}
\label{equation-filtration-cohomology}
F^pH^n(K^\bullet, d) =
\frac{\Ker(d) \cap F^pK^n + \Im(d) \cap K^n}{\Im(d) \cap K^n}
\end{equation}
and hence we see that
\begin{equation}
\label{equation-graded-cohomology}
\text{gr}^p H^n(K^\bullet) =
\frac{\Ker(d) \cap F^pK^n}{\Ker(d) \cap F^{p + 1}K^n + \Im(d) \cap F^pK^n}
\end{equation}
(one intermediate step omitted).

\begin{lemma}
\label{lemma-compute-cohomology-filtered-complex}
Let $\mathcal{A}$ be an abelian category. Let $(K^\bullet, F)$ be a filtered
complex of $\mathcal{A}$. If $Z_\infty^{p, q}$ and $B_\infty^{p, q}$ exist
(see proof), then
\begin{enumerate}
\item the limit $E_\infty$ exists and is a bigraded object having
$E_\infty^{p, q} = Z_\infty^{p, q}/B_\infty^{p, q}$ in bidegree $(p, q)$,
\item the $p$th graded part $\text{gr}^pH^n(K^\bullet)$ of the
$n$th cohomology object of $K^\bullet$ is a subquotient of
$E_\infty^{p, n - p}$.
\end{enumerate}
\end{lemma}

\begin{proof}
The objects $Z_\infty$, $B_\infty$, and the limit
$E_\infty = Z_\infty/B_\infty$ of
Definition \ref{definition-limit-spectral-sequence}
are bigraded objects of $\mathcal{A}$ by our construction of the
spectral sequence in Lemma \ref{lemma-spectral-sequence-filtered-complex}.
Since $Z_r^{p, q}$, resp.\ $B_r^{p, q}$ is the bidegree $(p, q)$
part of $Z_r$, resp.\ $B_r$, if we assume that
$$
Z_\infty^{p, q} = \bigcap\nolimits_r Z_r^{p, q} =
\bigcap\nolimits_r
\frac{F^pK^{p + q} \cap d^{-1}(F^{p + r}K^{p + q + 1}) + F^{p + 1}K^{p + q}}
{F^{p + 1}K^{p + q}}
$$
and
$$
B_\infty^{p, q} = \bigcup\nolimits_r B_r^{p, q} =
\bigcup\nolimits_r
\frac{F^pK^{p + q} \cap d(F^{p - r + 1}K^{p + q - 1}) + F^{p + 1}K^{p + q}}
{F^{p + 1}K^{p + q}}
$$
exist, then $Z_\infty$ and $B_\infty$ exist with bidegree $(p, q)$
parts $Z_\infty^{p, q}$ and $B_\infty^{p, q}$ (follows from an elementary
argument about unions and intersections of bigraded objects). Thus
$$
E_\infty^{p, q} =
\frac{\bigcap_r (F^pK^{p + q} \cap d^{-1}(F^{p + r}K^{p + q + 1})
+ F^{p + 1}K^{p + q})}
{\bigcup_r (F^pK^{p + q} \cap d(F^{p - r + 1}K^{p + q - 1})
+ F^{p + 1}K^{p + q})}.
$$
where the top and the bottom exist. With $n = p + q$ we have
\begin{equation}
\label{equation-on-top-bigraded}
\Ker(d) \cap F^pK^{n} + F^{p + 1}K^{n}
\subset
\bigcap\nolimits_r
\left(
F^pK^{n} \cap d^{-1}(F^{p + r}K^{n + 1}) + F^{p + 1}K^{n}
\right)
\end{equation}
and
\begin{equation}
\label{equation-at-bottom-bigraded}
\bigcup\nolimits_r
\left(
F^pK^{n} \cap d(F^{p - r + 1}K^{n - 1}) + F^{p + 1}K^{n}
\right)
\subset
\Im(d) \cap F^pK^{n} + F^{p + 1}K^{n}.
\end{equation}
Thus a subquotient of $E_\infty^{p, q}$ is
$$
\frac{\Ker(d) \cap F^pK^{n} + F^{p + 1}K^n}
{\Im(d) \cap F^pK^{n} + F^{p + 1}K^{n}} =
\frac{\Ker(d) \cap F^pK^n}{\Im(d) \cap F^pK^n + \Ker(d) \cap F^{p + 1}K^n}
$$
Comparing with (\ref{equation-graded-cohomology}) we conclude.
\end{proof}

\begin{definition}
\label{definition-bounded-ss}
Let $\mathcal{A}$ be an abelian category. Let $(E_r, d_r)_{r \geq r_0}$
be a spectral sequence of bigraded objects of
$\mathcal{A}$ with $d_r$ of bidegree $(r, -r + 1)$.
We say such a spectral sequence is
\begin{enumerate}
\item {\it regular} if for all $p, q \in \mathbf{Z}$ there is
a $b = b(p, q)$ such that the maps
$d_r^{p, q} : E_r^{p, q} \to E_r^{p + r, q - r + 1}$ are zero for $r \geq b$,
\item {\it coregular} if for all $p, q \in \mathbf{Z}$ there is a
$b = b(p, q)$ such that the maps
$d_r^{p - r, q + r - 1} : E_r^{p - r, q + r - 1} \to E_r^{p, q}$
are zero for $r \geq b$,
\item {\it bounded} if for all $n$
there are only a finite number of nonzero $E_{r_0}^{p, n - p}$,
\item {\it bounded below} if for all $n$ there is a $b = b(n)$ such that
$E_{r_0}^{p, n - p} = 0$ for $p \geq b$.
\item {\it bounded above} if for all $n$ there is a $b = b(n)$ such that
$E_{r_0}^{p, n - p} = 0$ for $p \leq b$.
\end{enumerate}
\end{definition}

\noindent
Bounded below means that if we look at $E_r^{p, q}$ on the line
$p + q = n$ (whose slope is $-1$) we obtain zeros as $(p, q)$
moves down and to the right.
As mentioned above there is no consistent terminology regarding
these notions in the literature.

\begin{lemma}
\label{lemma-relate-boundedness}
In the situation of Definition \ref{definition-bounded-ss}.
Let $Z_r^{p, q}, B_r^{p, q} \subset E_{r_0}^{p, q}$ be the
$(p, q)$-graded parts of $Z_r, B_r$ defined as in
Section \ref{section-spectral-sequence}.
\begin{enumerate}
\item The spectral sequence is regular if and only if for all $p, q$
there exists an $r = r(p, q)$ such that
$Z_r^{p, q} = Z_{r + 1}^{p, q} = \ldots$
\item The spectral sequence is coregular if and only if for all $p, q$
there exists an $r = r(p, q)$ such that
$B_r^{p, q} = B_{r + 1}^{p, q} = \ldots$
\item The spectral sequence is bounded if and only if it is both
bounded below and bounded above.
\item If the spectral sequence is bounded below, then it is regular.
\item If the spectral sequence is bounded above, then it is coregular.
\end{enumerate}
\end{lemma}

\begin{proof}
Omitted. Hint: If $E_r^{p, q} = 0$, then we have $E_{r'}^{p, q} = 0$
for all $r' \geq r$.
\end{proof}

\begin{definition}
\label{definition-filtered-complex-ss-converges}
Let $\mathcal{A}$ be an abelian category. Let $(K^\bullet, F)$ be a
filtered complex of $\mathcal{A}$. We say the spectral sequence
associated to $(K^\bullet, F)$
\begin{enumerate}
\item {\it weakly converges to $H^*(K^\bullet)$} if
$\text{gr}^pH^n(K^\bullet) = E_{\infty}^{p, n - p}$
via Lemma \ref{lemma-compute-cohomology-filtered-complex}
for all $p, n \in \mathbf{Z}$,
\item {\it abuts to $H^*(K^\bullet)$} if it weakly converges to
$H^*(K^\bullet)$ and $\bigcap_p F^pH^n(K^\bullet) = 0$ and
$\bigcup_p F^p H^n(K^\bullet) = H^n(K^\bullet)$ for all $n$,
\item {\it converges to $H^*(K^\bullet)$} if it is regular,
abuts to $H^*(K^\bullet)$, and
$H^n(K^\bullet) = \lim_p H^n(K^\bullet)/F^pH^n(K^\bullet)$.
\end{enumerate}
\end{definition}

\noindent
Weak convergence, abutment, or convergence is symbolized by the notation
$E_r^{p, q} \Rightarrow H^{p + q}(K^\bullet)$.
As mentioned above there is no consistent terminology regarding
these notions in the literature.

\begin{lemma}
\label{lemma-filtered-complex-ss-converges}
Let $\mathcal{A}$ be an abelian category. Let $(K^\bullet, F)$ be a filtered
complex of $\mathcal{A}$. The associated spectral sequence
\begin{enumerate}
\item weakly converges to $H^*(K^\bullet)$ if and only if for every
$p, q \in \mathbf{Z}$ we have equality in equations
(\ref{equation-at-bottom-bigraded}) and (\ref{equation-on-top-bigraded}),
\item abuts to $H^*(K)$ if and only if it weakly converges to $H^*(K^\bullet)$
and we have
$\bigcap_p (\Ker(d) \cap F^pK^n + \Im(d) \cap K^n) = \Im(d) \cap K^n$
and
$\bigcup_p (\Ker(d) \cap F^pK^n + \Im(d) \cap K^n) = \Ker(d) \cap K^n$.
\end{enumerate}
\end{lemma}

\begin{proof}
Immediate from the discussions above.
\end{proof}

\begin{lemma}
\label{lemma-biregular-ss-converges}
Let $\mathcal{A}$ be an abelian category. Let $(K^\bullet, F)$ be a
filtered complex of $\mathcal{A}$. Assume that the filtration on each $K^n$
is finite (see Definition \ref{definition-filtered}). Then
\begin{enumerate}
\item the spectral sequence associated to $(K^\bullet, F)$ is bounded,
\item the filtration on each $H^n(K^\bullet)$ is finite,
\item the spectral sequence associated to $(K^\bullet, F)$ converges
to $H^*(K^\bullet)$,
\item if $\mathcal{C} \subset \mathcal{A}$ is a weak Serre subcategory
and for some $r$ we have $E_r^{p, q} \in \mathcal{C}$ for all
$p, q \in \mathbf{Z}$, then $H^n(K^\bullet)$ is in $\mathcal{C}$.
\end{enumerate}
\end{lemma}

\begin{proof}
Part (1) follows as $E_0^{p, n - p} = \text{gr}^p K^n$.
Part (2) is clear from Equation (\ref{equation-filtration-cohomology}).
We will use Lemma \ref{lemma-filtered-complex-ss-converges} to prove
that the spectral sequence weakly converges. Fix $p, n \in \mathbf{Z}$.
The right hand side of (\ref{equation-on-top-bigraded})
is equal to $F^pK^n \cap \Ker(d) + F^{p + 1}K^n$ because
$F^{p + r}K^n = 0$ for $r \gg 0$. Thus (\ref{equation-on-top-bigraded})
is an equality. The left hand side of (\ref{equation-at-bottom-bigraded})
is equal to $F^pK^n \cap \Im(d) + F^{p + 1}K^n$ because
$F^{p - r + 1}K^{n - 1} = K^{n - 1}$ for $r \gg 0$.
Thus (\ref{equation-at-bottom-bigraded}) is an equality. Since the filtration
on $H^n(K^\bullet)$ is finite by (2) we see that we have abutment.
To prove we have convergence we have to show the spectral sequence is
regular which follows as it is bounded
(Lemma \ref{lemma-relate-boundedness}) and we have
to show that $H^n(K^\bullet) = \lim_p H^n(K^\bullet)/F^pH^n(K^\bullet)$
which follows from the fact that the filtration on $H^*(K^\bullet)$
is finite proved in part (2).

\medskip\noindent
Proof of (4). Assume that for some $r \geq 0$ we have
$E_r^{p, q} \in \mathcal{C}$ for some weak Serre subcategory
$\mathcal{C}$ of $\mathcal{A}$. Then $E_{r + 1}^{p, q}$ is
in $\mathcal{C}$ as well, see
Lemma \ref{lemma-characterize-weak-serre-subcategory}.
By boundedness proved above (which implies that the spectral sequence
is both regular and coregular, see Lemma \ref{lemma-relate-boundedness})
we can find an $r' \geq r$ such that $E_\infty^{p, q} = E_{r'}^{p, q}$
for all $p, q$ with $p + q = n$. Thus $H^n(K^\bullet)$ is an object
of $\mathcal{A}$ which has a finite filtration whose graded pieces
are in $\mathcal{C}$. This implies that $H^n(K^\bullet)$ is in $\mathcal{C}$
by Lemma \ref{lemma-characterize-weak-serre-subcategory}.
\end{proof}

\begin{lemma}
\label{lemma-biregular-ss-relation-in-K0}
Let $\mathcal{A}$ be an abelian category. Let $(K^\bullet, F)$ be a
filtered complex of $\mathcal{A}$. Assume that the filtration on each $K^n$
is finite (see Definition \ref{definition-filtered}) and that for some
$r$ we have only a finite number of nonzero $E_r^{p, q}$. Then
only a finite number of $H^n(K^\bullet)$ are nonzero and we have
$$
\sum (-1)^n[H^n(K^\bullet)] = \sum (-1)^{p + q} [E_r^{p, q}]
$$
in $K_0(\mathcal{A}')$ where $\mathcal{A}'$ is the smallest weak
Serre subcategory of $\mathcal{A}$ containing the objects
$E_r^{p, q}$.
\end{lemma}

\begin{proof}
Denote $E_r^{even}$ and $E_r^{odd}$ the even and odd part of $E_r$
defined as the direct sum of the $(p, q)$ components with $p + q$ even
and odd. The differential $d_r$ defines maps
$\varphi : E_r^{even} \to E_r^{odd}$ and $\psi : E_r^{odd} \to E_r^{even}$
whose compositions either way give zero.
Then we see that
\begin{align*}
[E_r^{even}] - [E_r^{odd}] & =
[\Ker(\varphi)] + [\Im(\varphi)] - [\Ker(\psi)] - [\Im(\psi)] \\
& =
[\Ker(\varphi)/\Im(\psi)] - [\Ker(\psi)/\Im(\varphi)] \\
& =
[E_{r + 1}^{even}] - [E_{r + 1}^{odd}]
\end{align*}
Note that all the intervening objects are in the smallest Serre
subcategory containing the objects $E_r^{p, q}$.
Continuing in this manner we see that we can increase $r$ at will.
Since there are only a finite number of pairs $(p, q)$ for which
$E_r^{p, q}$ is nonzero, a property which is inherited by
$E_{r + 1}, E_{r + 2}, \ldots$, we see that we may assume that $d_r = 0$.
At this stage we see that $H^n(K^\bullet)$ has a finite filtration
(Lemma \ref{lemma-biregular-ss-converges}) whose graded pieces
are exactly the $E_r^{p, n - p}$ and the result is clear.
\end{proof}

\noindent
The following lemma is more a kind of sanity check for our definitions.
Surely, if we have a filtered complex such that for every $n$ we have
$$
H^n(F^pK^\bullet) = 0\text{ for }p \gg 0
\quad\text{and}\quad
H^n(F^pK^\bullet) = H^n(K^\bullet)\text{ for }p \ll 0,
$$
then the corresponding spectral sequence should converge?

\begin{lemma}
\label{lemma-ss-converges-trivial}
Let $\mathcal{A}$ be an abelian category. Let $(K^\bullet, F)$ be a
filtered complex of $\mathcal{A}$. Assume
\begin{enumerate}
\item for every $n$ there exist $p_0(n)$ such that
$H^n(F^pK^\bullet) = 0$ for $p \geq p_0(n)$,
\item for every $n$ there exist $p_1(n)$ such that
$H^n(F^pK^\bullet) \to H^n(K^\bullet)$ is an isomorphism
for $p \leq p_1(n)$.
\end{enumerate}
Then
\begin{enumerate}
\item the spectral sequence associated to $(K^\bullet, F)$ is bounded,
\item the filtration on each $H^n(K^\bullet)$ is finite,
\item the spectral sequence associated to $(K^\bullet, F)$ converges
to $H^*(K^\bullet)$.
\end{enumerate}
\end{lemma}

\begin{proof}
Fix $n$. Using the long exact cohomology sequence associated to
the short exact sequence of complexes
$$
0 \to F^{p + 1}K^\bullet \to F^pK^\bullet \to \text{gr}^pK^\bullet \to 0
$$
we find that $E_1^{p, n - p} = 0$ for $p \geq \max(p_0(n), p_0(n + 1))$ and
$p < \min(p_1(n), p_1(n + 1))$. Hence the spectral sequence is bounded
(Definition \ref{definition-bounded-ss}). This proves (1).

\medskip\noindent
It is clear from the assumptions and
Definition \ref{definition-filtration-cohomology-filtered-complex}
that the filtration on $H^n(K^\bullet)$ is finite. This proves (2).

\medskip\noindent
Next we prove that the spectral sequence weakly converges to
$H^*(K^\bullet)$ using
Lemma \ref{lemma-filtered-complex-ss-converges}.
Let us show that we have equality in (\ref{equation-on-top-bigraded}).
Namely, for $p + r > p_0(n + 1)$ the map
$$
d : F^pK^{n} \cap d^{-1}(F^{p + r}K^{n + 1}) \to F^{p + r}K^{n + 1}
$$
ends up in the image of $d : F^{p + r}K^n \to F^{p + r}K^{n + 1}$
because the complex $F^{p + r}K^\bullet$ is exact in degree $n + 1$.
We conclude that $F^pK^{n} \cap d^{-1}(F^{p + r}K^{n + 1}) =
d(F^{p + r}K^n) + \Ker(d) \cap F^pK^n$. Hence for such $r$ we have
$$
\Ker(d) \cap F^pK^{n} + F^{p + 1}K^{n} =
F^pK^{n} \cap d^{-1}(F^{p + r}K^{n + 1}) + F^{p + 1}K^{n}
$$
which proves the desired equality. To show that we have equality in
(\ref{equation-at-bottom-bigraded}) we use that for $p - r + 1 < p_1(n - 1)$
we have
$$
d(F^{p - r + 1}K^{n - 1}) = \Im(d) \cap F^{p - r + 1}K^n
$$
because the map $F^{p - r + 1}K^\bullet \to K^\bullet$ induces an
isomorphism on cohomology in degree $n - 1$. This shows that
we have
$$
F^pK^{n} \cap d(F^{p - r + 1}K^{n - 1}) + F^{p + 1}K^{n} =
\Im(d) \cap F^pK^{n} + F^{p + 1}K^{n}
$$
for such $r$ which proves the desired equality.

\medskip\noindent
To see that the spectral sequence abuts to $H^*(K^\bullet)$ using
Lemma \ref{lemma-filtered-complex-ss-converges} we have to show that
$\bigcap_p (\Ker(d) \cap F^pK^n + \Im(d) \cap K^n) = \Im(d) \cap K^n$
and
$\bigcup_p (\Ker(d) \cap F^pK^n + \Im(d) \cap K^n) = \Ker(d) \cap K^n$.
For $p \geq p_0(n)$ we have
$\Ker(d) \cap F^pK^n + \Im(d) \cap K^n = \Im(d) \cap K^n$
and for $p \leq p_1(n)$ we have
$\Ker(d) \cap F^pK^n + \Im(d) \cap K^n = \Ker(d) \cap K^n$.
Combining weak convergence, abutment, and boundedness we see
that (2) and (3) are true.
\end{proof}






















